\section{Results}
\label{sec:results}

Table~\ref{tab:main} summarises the study; Figure~\ref{fig:support} plots
selection quality and active-set size against $E$; the complete grid is in
Table~\ref{tab:full}. All entries are means over $200$ replicates per
design, with Monte Carlo standard errors.

\begin{table}[t]
\centering
\small
\caption{Monte Carlo results over 100 replicates per design ($p=600$, $s=30$, $K=3$, $m=(400,240,160)$). Support metrics are against the true support at $\epsilon=10^{-4}$. Monte Carlo standard errors for $F_1$ in parentheses. \textsc{FedAvg} and \textsc{FedAvg-ST} are shown at $E=1$ and $E=20$; the full grid is in Table~\ref{tab:full}.}
\label{tab:main}
\begin{tabular}{llrrrrrrr}
\toprule
Design & Method & Active & Prec. & Recall & $F_1$ & $\|\hat\beta-\beta^\star\|_2$ & Test MSE & Rounds \\
\midrule
Independent & Pooled & 127.6 & 0.246 & 0.977 & 0.388 (0.008) & 1.94 & 24.43 & --- \\
 & Local-only & 105.5 & 0.260 & 0.863 & 0.395 (0.007) & 2.85 & 28.81 & 0.0 \\
 & One-shot & 198.4 & 0.150 & 0.940 & 0.257 (0.006) & 2.91 & 29.18 & 1.0 \\
 & ADMM & 43.6 & 0.681 & 0.875 & 0.738 (0.011) & 2.51 & 27.19 & 58.7 \\
 & FedAvg ($E=1$) & 253.6 & 0.125 & 0.955 & 0.218 (0.006) & 2.85 & 28.82 & 14.6 \\
 & FedAvg ($E=20$) & 203.1 & 0.149 & 0.940 & 0.255 (0.006) & 2.91 & 29.17 & 3.6 \\
 & FedAvg-ST ($E=1$) & 45.4 & 0.576 & 0.801 & 0.654 (0.009) & 3.21 & 31.03 & 15.6 \\
 & FedAvg-ST ($E=20$) & 46.0 & 0.565 & 0.795 & 0.644 (0.009) & 3.27 & 31.41 & 4.6 \\
\midrule
Correlated & Pooled & 121.1 & 0.255 & 0.968 & 0.399 (0.008) & 2.09 & 24.61 & --- \\
 & Local-only & 107.7 & 0.252 & 0.846 & 0.383 (0.007) & 3.03 & 28.91 & 0.0 \\
 & One-shot & 198.8 & 0.146 & 0.930 & 0.251 (0.005) & 3.04 & 28.83 & 1.0 \\
 & ADMM & 48.3 & 0.601 & 0.860 & 0.681 (0.011) & 2.58 & 27.15 & 93.8 \\
 & FedAvg ($E=1$) & 263.2 & 0.117 & 0.949 & 0.206 (0.005) & 2.96 & 28.42 & 16.5 \\
 & FedAvg ($E=20$) & 201.6 & 0.145 & 0.930 & 0.250 (0.005) & 3.03 & 28.83 & 3.8 \\
 & FedAvg-ST ($E=1$) & 59.3 & 0.457 & 0.805 & 0.565 (0.010) & 3.22 & 30.70 & 17.5 \\
 & FedAvg-ST ($E=20$) & 61.4 & 0.428 & 0.796 & 0.543 (0.010) & 3.29 & 31.08 & 4.8 \\
\midrule
Heterogeneous & Pooled & 116.1 & 0.256 & 0.888 & 0.388 (0.009) & 2.77 & 38.55 & --- \\
 & Local-only & 105.5 & 0.260 & 0.863 & 0.395 (0.007) & 2.85 & 39.23 & 0.0 \\
 & One-shot & 185.3 & 0.157 & 0.922 & 0.266 (0.005) & 3.10 & 40.25 & 1.0 \\
 & ADMM & 44.0 & 0.603 & 0.646 & 0.537 (0.014) & 3.33 & 42.28 & 66.1 \\
 & FedAvg ($E=1$) & 258.5 & 0.121 & 0.943 & 0.212 (0.005) & 3.05 & 40.05 & 14.2 \\
 & FedAvg ($E=20$) & 194.0 & 0.154 & 0.923 & 0.262 (0.006) & 3.10 & 40.24 & 4.0 \\
 & FedAvg-ST ($E=1$) & 23.5 & 0.830 & 0.593 & 0.654 (0.016) & 3.72 & 44.93 & 15.2 \\
 & FedAvg-ST ($E=20$) & 25.4 & 0.807 & 0.605 & 0.652 (0.013) & 3.75 & 45.14 & 5.0 \\

\bottomrule
\end{tabular}
\end{table}

\begin{figure}[t]
\centering
\includegraphics[width=\textwidth]{figures/fig1_support.pdf}
\caption{Selection quality (top) and active-set size (bottom) against the
number of local epochs $E$, by design. Error bars are Monte Carlo standard
errors over $200$ replicates. Horizontal lines are comparators that do not
depend on $E$, shown under the one-standard-error tuning rule so that all
curves and lines are selected on the same footing. The true model has
$s = 30$ non-zero coefficients out of $p = 600$.}
\label{fig:support}
\end{figure}

\subsection{Plain averaging returns a dense estimator at every $E$}

The first row of Figure~\ref{fig:support} is dominated by a gap that has
nothing to do with $E$. With $s = 30$ true signals out of $p = 600$,
\textsc{FedAvg} declares \MIndFedAct, \MCorFedAct\ and \MHetFedAct\
coefficients active at $E = 1$ in the independent, correlated and
heterogeneous designs---between a third and a half of all predictors. The
signature is the one predicted by Proposition~\ref{prop:orth}(iv): recall is
very high (\MIndFedRec, \MCorFedRec, \MHetFedRec) and precision very low
(\MIndFedPrec, \MCorFedPrec, \MHetFedPrec). The aggregated estimator misses
almost no true signal and admits almost everything else, which is what a
union of three local active sets looks like.

Increasing $E$ shrinks the active set---by \MIndActDrop, \MCorActDrop\ and
\MHetActDrop\ from $E = 1$ to $E = 20$---and correspondingly raises $F_1$
from \MIndFedF\ to \MIndFedTF, \MCorFedF\ to \MCorFedTF, and \MHetFedF\ to
\MHetFedTF. It would be natural to read this as more local work producing
better selection. The reading we prefer is that more local work produces
\emph{less union}: each site's iterate is closer to its own sparse local
minimiser, so fewer distinct coordinates survive the average. Consistently
with Proposition~\ref{prop:orth}(ii), the shrinkage is smallest in the
independent design and largest in the heterogeneous one, though the ordering
is modest and we would not lean on it: with $p > m_j$ no design here is close
to the orthogonal case in which the $E$-effect is exactly zero.

The practical conclusion is the same at both ends of the range. At every $E$
and in every design, plain averaging recovers the true support less well than
\emph{doing nothing}: the Lasso at the largest site alone gives
$F_1 = \MIndLocSeF$, \MCorLocSeF\ and \MHetLocSeF\ under matched tuning,
against at most \MIndFedTF, \MCorFedTF\ and \MHetFedTF\ for \textsc{FedAvg}.
Nor does the communication buy much: by $E = 20$, \textsc{FedAvg} has
converged to essentially the one-shot average, which uses a single round
($F_1 = \MIndOneF$, \MCorOneF, \MHetOneF).

\subsection{The tuning rule moves selection as much as the aggregation rule}

Before crediting any repair of the aggregation with the improvement it
appears to produce, it is worth seeing how much of the same improvement is
available from the penalty-selection rule alone. Compare the two centralised
rows in Table~\ref{tab:main}. Fitting the pooled Lasso with the penalty that
minimises validation MSE gives $F_1 = \MIndPoolF$, \MCorPoolF\ and
\MHetPoolF; fitting exactly the same estimator on exactly the same data with
the one-standard-error rule gives \MIndPoolSeF, \MCorPoolSeF\ and
\MHetPoolSeF. Changing nothing but the rule moves $F_1$ by more than the
entire range we observe across $E$.

The same holds for the exact federated solver. Consensus ADMM inherits the
penalty $\bar\lambda$ from the sites' local choices; selecting the consensus
penalty by federated validation instead gives \textsc{ADMM-CV}, whose $F_1$
collapses to \MIndAdmmCvF, \MCorAdmmCvF\ and \MHetAdmmCvF\ under the
minimum-MSE rule and recovers to \MIndAdmmCvSeF, \MCorAdmmCvSeF\ and
\MHetAdmmCvSeF\ under the one-standard-error rule. Exact optimisation of the
consensus objective is therefore worth much less than it first appears: most
of what \textsc{ADMM} achieves at $\bar\lambda$ is attributable to
$\bar\lambda$ being, in these designs, a sparsity-favouring penalty.

We draw two conclusions. First, comparisons of federated selection procedures
must fix the selection rule across all arms; we report every comparator under
both rules for that reason, and read the results down the \textsc{-1SE} rows.
Second, and more substantively, a study that varies $E$ while holding a
prediction-tuned penalty fixed is varying the wrong knob: the penalty rule
dominates.

\subsection{Repairing the aggregation: thresholding after, or during}

Against matched baselines, both repairs work and neither is dramatic.
\textsc{FedAvg-ST} reaches $F_1 = \MIndStF$ (SE \MIndStSe), \MCorStF\
(\MCorStSe) and \MHetStF\ (\MHetStSe) at $E = 1$, roughly three times the
plain-averaging value, with active sets of \MIndStAct, \MCorStAct\ and
\MHetStAct. Per-round proximal aggregation, \textsc{FedAvg-P}, is close
behind at \MIndProxF, \MCorProxF\ and \MHetProxF, rising slightly with $E$ to
\MIndProxTF, \MCorProxTF\ and \MHetProxTF. That the cheap retrofit matches
the in-loop correction is worth knowing, since the retrofit needs no change
to the training loop; we would not read the ordering between them as stable.

Both remain below a well-tuned centralised fit in the two homogeneous designs
(\MIndPoolSeF\ and \MCorPoolSeF\ for \textsc{Pooled-1SE}), which is as it
should be---no federated procedure should beat the pooled estimator it
approximates. In the heterogeneous design \textsc{FedAvg-ST} (\MHetStF)
edges above it (\MHetPoolSeF), which we read as the validation-tuned
threshold adapting to site-level scale differences that a single pooled
penalty cannot.

The improvement is not free, and the direction of the trade matters.
Thresholding removes small non-zero coefficients that were contributing to
fit, so estimation and prediction deteriorate: in the independent design
$\|\hat\bbeta - \bbeta^\star\|_2$ rises from \MIndFedErr\ to \MIndStErr\ and
test MSE from \MIndFedMse\ to \MIndStMse. An analyst who wants a variable
list should take that trade; one who wants a predictor should not, and should
use the prediction-tuned pooled or \textsc{ADMM-CV} fit
(\MIndPoolMse\ and \MIndAdmmCvMse\ test MSE).

Finally, and unlike plain averaging, the thresholded schedules are close to
flat in $E$: \MIndStF\ to \MIndStTF, \MCorStF\ to \MCorStTF\ and \MHetStF\ to
\MHetStTF\ across the whole grid, changes of one to two Monte Carlo standard
errors. In these designs, once the aggregation rule is repaired, $E$ can be
chosen on cost rather than on its effect on the variable list. That statement
is design-specific rather than general: Section~\ref{sec:realdata} reports a
data set on which the flattening is only partial.

\subsection{Exact consensus is not a universal answer}

Consensus ADMM at $\bar\lambda$ gives the best support recovery of any method
we consider in the independent (\MIndAdmmF) and correlated (\MCorAdmmF)
designs, with \MIndAdmmAct\ and \MCorAdmmAct\ active coefficients. Holding
the data, the objective, the penalty and the privacy boundary fixed and
changing only how the sites are combined moves $F_1$ from \MIndFedF\ to
\MIndAdmmF---the cleanest statement of the article's first point.

The heterogeneous design is the exception, and it is instructive. There ADMM
attains only \MHetAdmmF, below \textsc{FedAvg-ST} (\MHetStF), below
\textsc{FedAvg-P} (\MHetProxTF\ at $E = 20$) and below the local-only
baseline (\MHetLocSeF). Exactness is a property of the optimisation, not of
the estimand: ADMM solves the pooled Lasso at a single common penalty, and
when sites differ in design scale and noise level no single penalty serves
them all. Exact consensus is also the most expensive option here, needing
\MIndAdmmRounds, \MCorAdmmRounds\ and \MHetAdmmRounds\ rounds against
\MIndRoundsOne, \MCorRoundsOne\ and \MHetRoundsOne\ for \textsc{FedAvg} at
$E = 1$.

\subsection{Communication and computation move in opposite directions}

\begin{figure}[t]
\centering
\includegraphics[width=\textwidth]{figures/fig2_cost.pdf}
\caption{The two cost axes for \textsc{FedAvg}. Upload (solid, left axis)
falls in $E$; total coordinate-descent epochs per site (dashed, right axis)
rises. Means over $200$ replicates.}
\label{fig:cost}
\end{figure}

Figure~\ref{fig:cost} makes the accounting point. In the independent design,
moving from $E = 1$ to $E = 20$ cuts rounds from \MIndRoundsOne\ to
\MIndRoundsTwenty\ and upload in proportion, but raises total local work from
\MIndEpochsOne\ to \MIndEpochsTwenty\ coordinate-descent epochs per site---a
factor of five. The correlated and heterogeneous designs behave the same way
(\MCorEpochsOne\ to \MCorEpochsTwenty, and \MHetEpochsOne\ to
\MHetEpochsTwenty). A schedule selected on rounds alone performs the most
arithmetic, and in cross-silo deployments, where sites are institutions with
real compute budgets rather than phones, that is not obviously the right
optimum.

The threshold-selection round is negligible on this scale. It transmits
\MIndStScalar~KiB of scalars in total, against \MIndFedVectorOne~KiB of
coefficient vectors for a single \textsc{FedAvg} round at $E = 1$; the two
are reported separately in Table~\ref{tab:full}.
